% This file is part of proScientia.ltx
% (c) 2022 Karsten Reincke (https://github.com/kreincke/proScientia.ltx)
% It is distributed under the terms of the creative commons license
% CC-BY-4.0 (= https://creativecommons.org/licenses/by/4.0/)

\documentclass[
  DIV=calc,
  BCOR=5mm,
  11pt,
  headings=small,
  oneside,
  abstract=true,
  toc=bib,
  english,ngerman]{scrartcl}

%%% (1) general configurations %%%
\usepackage[utf8]{inputenc}
\usepackage{a4}
\usepackage[ngerman]{babel}

\usepackage[
  backend=biber,
  style=authortitle-dw,
  sortlocale=auto,
]{biblatex}
\ProvidesFile{biblatex.cfg}
% This file is part of proScientia.ltx
% (c) 2022 Karsten Reincke (https://github.com/kreincke/proScientia.ltx)
% It is distributed under the terms of the creative commons license
% CC-BY-4.0 (= https://creativecommons.org/licenses/by/4.0/)


\ExecuteBibliographyOptions{
  %%%%%%%%%%%%%%%%%%%%%%%%%%%%%%%%%%%%%%%%%%%%
  % configurations offered by biblatex itself
  % ------------------------------------------
  maxnames=5,  % Truncate author list after 5 authors ...
  minnames=3,  % ... But display at least 3 authors
  autocite=inline,
  hyperref=true,  % Use hyperref package (should be automatically detected, though)
  backref=true,  % Back references from bibliography page to each encounter
  backrefstyle=two,  % Combine back refs if on two consecutive pages
  isbn=true,  % (Dont) print ISBN, ISSN numbers
  autolang=hyphen,
  % track 'reused' csquotes
  citetracker=constrict, %
  loccittracker=constrict, % discriminate different pages (a.a.O) versues same page (ebda)
  opcittracker=constrict,
  idemtracker=constrict,
  ibidtracker=constrict,
  pagetracker=true,
  %%%%%%%%%%%%%%%%%%%%%%%%%%%%%%%%%%%%%%%%%%%%
  % configurations offered by authortitle-dw
  % ------------------------------------------
  annotation =true,
  namefont=normal,
  firstnamefont=normal,
  idemfont=italic,
  ibidemfont=italic,
  idembib=true, % cluster the books of the same author in bib
  idembibformat=idem, % indicate the same author by ders/dies.
  editorstring=brackets, % parens=(Hrsg.) | brackets=[Hrsg.] | normal = , Hrsg.
  nopublisher=false, % insert publisher into bib data
  editionstring=true, % allow strings in the edition field
  % ZNAME VOL (YEAR) Nr. YOURNALNUMBER
  journalnumber=afteryear, %
  %journumstring=h.
  series=afteryear,
  seriesformat=parens,
  addyear=true, %insert year after titel in 'shorttitle' hints => no year in shorttitle
  firstfull=true, % frist quote complete
  edstringincitations=false, %editor and translator only in the first
  citepages=separate, % vollzitat erst mit seiten, dann 'hier: S. 12'
}
% Put your definitions here.

% refine seriesformat by adding a prefix inside of the parens
\renewcommand*{\seriespunct}{=\addspace}

% by default Biblatex-dw uses Autor1/Autor2 in cites
% these redefinitions overwrite that behaviour by
% duplicating the style used for the bibliography
% (S. p. 35 in the German biblatex-dw handbook )
\renewcommand*{\citemultinamedelim}{\addcomma\space}
\renewcommand*{\citefinalnamedelim}{%
\ifnum\value{liststop}>2 \finalandcomma\fi
\addspace\bibstring{and}\space}%
\renewcommand*{\citerevsdnamedelim}{\addspace}

% biblatex-dw does not print 'ders.' / 'dies.' in a row of cites quoting
% the same book. Additionally, it does not know the diffrence between the
% same page of of the same work and a different page of the same work
% quoted before.
%
% As soon as biblatex 3.15 is offered by UBUNTU
% use \bibncpstring[\mkibid]{ibidem} instaed of inserting the
% German string literally as this hack does:
\renewbibmacro*{cite:ibid}{%
{\ifthenelse{\ifloccit}
  {\printtext[bibhyperref]{\usebibmacro{cite:idem} \mkibid{ebda}}%
   \global\booltrue{cbx:loccit}}
  {\printtext[bibhyperref]{\usebibmacro{cite:idem} \mkibid{a.a.O, }}}
}
}%
\DefineBibliographyStrings{english}{%
  urlseen = {RDL},
}
\DefineBibliographyStrings{german}{%
  urlseen = {RDL},
}
\endinput


\addbibresource{bib/lit.verify.bib}
\addbibresource{bib/lit.main.bib}

% package for improving the grey value and the line feed handling
\usepackage{microtype}

%language specific quoting signs
\usepackage[
  style=german,
  autostyle=true,
]{csquotes}

% language specific hyphenation
% This file is part of proScientia.ltx
% (c) 2022 Karsten Reincke (https://github.com/kreincke/proScientia.ltx)
% It is distributed under the terms of the creative commons license
% CC-BY-4.0 (= https://creativecommons.org/licenses/by/4.0/)

\hyphenation{ pro-scien-tia there-fo-re}

\babelhyphenation{ pro-scien-tia wor-tmi-tfalsch-entr-ennu-ngen}


%%% (3) layout page configuration %%%

% select the visible parts of a page
% S.31: { plain|empty|headings|myheadings }
\pagestyle{headings}

% select the wished style of page-numbering
% S.32: { arabic,roman,Roman,alph,Alph }
\pagenumbering{arabic}
\setcounter{page}{1}

% select the wished distances using the general setlength order:
% S.34 { baselineskip| parskip | parindent }
% - general no indent for paragraphs
\setlength{\parindent}{0pt}
\setlength{\parskip}{1.2ex plus 0.2ex minus 0.2ex}


%- start(footnote-configuration)

\deffootnote[1.5em]{1.5em}{1.5em}{\textsuperscript{\thefootnotemark)\ }}

%for using label as nameref
\usepackage{nameref}

%integrate nomenclature
% This file is part of proScientia.ltx
% (c) 2022 Karsten Reincke (https://github.com/kreincke/proScientia.ltx)
% It is distributed under the terms of the creative commons license
% CC-BY-4.0 (= https://creativecommons.org/licenses/by/4.0/)

\usepackage[intoc]{nomencl}
\let\abbr\nomenclature
% Deutsche Überschrift
\renewcommand{\nomname}{Abkürzungen}

\setlength{\nomlabelwidth}{.20\hsize}
\renewcommand{\nomlabel}[1]{#1 \dotfill}
% reduce the line distance
\setlength{\nomitemsep}{-\parsep}
\makenomenclature


% Hyperlinks
\usepackage{hyperref}
\hypersetup{bookmarks=true,breaklinks=true,colorlinks=true,citecolor=blue,draft=false}

\usepackage[T1]{fontenc}
\usepackage{wesu}
\usepackage{suetterl} % provided by fundus-suetterlin of TexLive

\usepackage{fontspec}
\newfontfamily\suetterhjz{Suetterlin HJZ 1911 Italic2024-03}

\usepackage{fontawesome6}

\newcommand{\sz}{{\ss}}

\begin{document}

%% use all entries of the bliography
\nocite{*}

\titlehead{proSuetterlin.ltx}
\subject{Release 0.1}
\title{Sütterlin in \LaTeX-Texten}
\subtitle{Ein selbsterklärendes Demo / Template}
\author{Karsten Reincke% This file is part of proScientia.ltx
% (c) 2022 Karsten Reincke (https://github.com/kreincke/proScientia.ltx)
% It is distributed under the terms of the creative commons license
% CC-BY-4.0 (= https://creativecommons.org/licenses/by/4.0/)

\footnote{\textbf{Dieser Text wird unter der XYZ Lizenz veröffentlicht.}
Hier können Ihre Bedingungen stehen, unter denen Sie Ihren eigenen Text weitergeben.
Gute Kandidaten wären z.B. die Creative Commons Lizenzen
\href{https://creativecommons.org/}{https://creativecommons.org/}. Traditionell ist auch die Formel: \emph{Alle Rechte vorbehalten. Die Verwendung von Text und Bildern, auch auszugsweise, bedarf der schriftlichen Zustimmung}. \newline
Da Ihre Arbeit auf dem Templatesystem \textit{proSuetterliv.ltx} aufbaut und da Sie dieses unter den Bedingungen der \texttt{CC-BY-4.0} Lizenz erhalten haben, müssen Sie auf dessen Verwendung hinweisen. Eine lizenzerfüllende Notiz könnte sein:
\newline
\textit{Erstellt auf der Basis des CC-BY-4.0 lizenzierten Tools \texttt{proSuetterlin.ltx} von K. Reincke \copyright{} 2026 [
Repository \href{https://github.com/kreincke/proSuetterlin.ltx}{https://github.com/kreincke/proSuetterlin.ltx} ,
Lizenztext \href{https://creativecommons.org/licenses/by/4.0/}{https://creativecommons.org/licenses/by/4.0/} ]
}}
}

\maketitle

\section{Sütterlin unter \LaTeX}

Bisher habe ich drei Methoden gefunden, Sütterlin in einem \LaTeX-Text zu verwenden:

\begin{enumerate}
  \item \texttt{\textbackslash{}usepackage[T1]\{fontenc\}} und \texttt{\textbackslash{}usepackage\{wesu\}} im Header, sowie \texttt{\{\textbackslash{}wesu ... \}} im Text
  \item \texttt{\textbackslash{}usepackage[T1]\{fontenc\}}, \texttt{\textbackslash{}usepackage\{suetterl\}} im Header, sowie \texttt{\{\textbackslash{}suetterlin ... \}} im Text
  \item \texttt{\textbackslash{}usepackage\{fontspec\}}, \texttt{\textbackslash{}newfontfamily\textbackslash{}suetterhjz\{Suetterlin HJZ 1911 Italic2024-03\}} im Header, sowie \texttt{\{\textbackslash{}suetterhjz ...\}} im Text, wobei man zuvor die Schrift von \href{https://zinken.net/Fonts/Suetterlin.html}{https://zinken.net/Fonts/Suetterlin.html} herunterladen, unter \texttt{.local/share/fonts/} als TrueType-Font ablegen und mittels \texttt{sudo fs-cache} bekanntmachen muss.
\end{enumerate}

Laut Nöhrenberg gibt es 7 Fälle, in denen in Sütterlintexten das \textit{Runde-s} anstelle des \textit{Langen-s} verwendet werden muss.\footcite[vgl.][15ff]{Noehrenberg2022b} Leider vermag keine der hier dargestellten Varianten diese Regeln aus sich heraus komplett abzudecken. 

Allerdings kann man beim Font \textit{wesu} die Verwendung des \textit{Runden-s} erzwingen, indem man die Silbentrennung explizit mit \texttt{\textbackslash -} in das geschriebene Wort einbaut. Dann darf man in den anderen Fällen aber nicht manuell in die Silbentrennung eingreifen. Und man muss akzeptieren, dass - im Quelltext geschrieben: \textit{Kios\textbackslash{}-k} - als {\wesu Kios\-k} erscheint und im Falle eines Zeilenumbruchs irregulär zwischen \texttt{s} und \texttt{k} getrennt wird.

Außerdem gibt es laut Nöhrenberg in Sütterlin ein gesondertes Zeichen für das \texttt{ß}, nämlich {\wesu {\ss}}. Dies ist aber offensichtlich nur über die alte \LaTeX-Notation \texttt{\textbackslash ss} erreichbar. Das 8-Bit-Ascii- bzw. UTF-8-Zeichen \texttt{ß} literal im \LaTeX-Text zu verwenden, hilft nicht.

Damit das angrenzende Leerzeichen nicht unterdrückt wird, muss das Syntagma für das Zeichen {\wesu {\ss}} im \LaTeX-Quelltext für den Sütterlintext noch in geschweifte Klammern gesetzt werden, was das etwas umständliche Syntagma \texttt{\{\textbackslash ss\}} ergäbe. Zur Vereinfachung habe ich einen entsprechenden Befehl in der Präambel definiert:

\begin{verbatim}
  \newcommand{\sz}{{\ss}}
\end{verbatim}

So bleibt bisher als Fazit, dass zwar die Schrift \textit{Suetterlin-HJZ-Italic\_1911.ttf} das vielleicht schönere Schriftbild bietet, dass aber nur mit \textit{wesu} die Sütterlinregeln in/mit {\LaTeX} umgesetzt werden können.

\subsection{Writing \textit{Sütterlin} with package \textit{wesu}}

{\wesu Hier ein Text, geschrieben in Sütterlin via \texttt{xelatex} und Package \texttt{wesu}}

Hier die Schluss-S-Probe:

\begin{itemize}
  \item \textit{Rundes-S}:
  \begin{itemize}
    \item[\faCheck] am Wortende: Haus = {\wesu Haus}, Maus = {\wesu Maus}, Tennis = {\wesu Tennis} 
    \item[\faCheck] in zusammengesetzten Wörtern: Marsmensch = {\wesu Mars\-mensch}, Haustür = {\wesu Haus\-tür} 
    \item[\faCheck] bei Fugenmorphem: Königstochter = {\wesu Königs\-tochter}, Valentinstag ={\wesu Valentins\-tag} 
    \item[\faCheck] vor konsonantischen Nachsilben: böse = {\wesu böse} versus Bos\-heit = {\wesu Bosheit} 
    \item[\faCheck] am Ende der Vorsilben 'dis' und 'des': Desinformation = {\wesu Des\-information}, Disharmonie = {\wesu Dis\-harmonie} 
    \item[\faCheck] am Wortende vor k: Kiosk = {\wesu Kios\-k} grotesk = {\wesu grotes\-k} 
    \item[\faCheck] am Silbenende vor d, k, m, n, w: Potsdam = {\wesu Pots\-dam}, Charisma= {\wesu Charis\-ma}
  \end{itemize}
  \item[\faCheck] \textit{Langes s} im Wort und sonst: Mensch = {\wesu Mensch}, hasten = {\wesu hasten}, Tasse = {\wesu Tasse} 
  \item[\faCheck] ß wie früher üblich: Haß = {\wesu Ha\ss} , daß = {\wesu da\ss} 
\end{itemize}

\subsection{Writing \textit{Sütterlin} with package \textit{suetterl}}

{\suetterlin Hier ein Text, geschrieben in Sütterlin via \texttt{xelatex} und Package \texttt{sutterl}, enthalten in CTAN-Paket \texttt{fundus-suetterlin}}

\begin{itemize}
  \item \textit{Rundes-S}:
  \begin{itemize}
    \item[\faXmark] am Wortende: Haus = {\suetterlin Haus}, Maus = {\suetterlin Maus}, Tennis = {\suetterlin Tennis} 
    \item[\faXmark] in zusammengesetzten Wörtern: Marsmensch = {\suetterlin Mars\-mensch}, Haustür = {\suetterlin Haus\-tür}
    \item[\faXmark] bei Fugenmorphem: Königstochter = {\suetterlin Königs\-tochter}, Valentinstag ={\suetterlin Valentins\-tag}
    \item[\faXmark] vor konsonantischen Nachsilben: böse = {\suetterlin böse} versus Bosheit = {\suetterlin Bos\-heit}
    \item[\faXmark] am Ende der Vorsilben 'dis' und 'des': Desinformation = {\suetterlin Des\-information}, Disharmonie = {\suetterlin Dis\-harmonie}
    \item[\faXmark] am Wortende vor k: Kiosk = {\suetterlin Kios\-k} grotesk = {\suetterlin grotes\-k}
    \item[\faXmark] am Silbenende vor d, k, m, n, w: Potsdam = {\suetterlin Pots\-dam}, Charisma= {\suetterlin Charis\-ma}
  \end{itemize}
  \item[\faCheck] \textit{Langes s} im Wort und sonst: Mensch = {\suetterlin Mensch}, hasten = {\suetterlin hasten}, Tasse = {\suetterlin Tasse}
  \item[\faCheck] ß wie früher üblich: Haß = {\suetterlin Ha\ss} , daß = {\wesu da\ss} 
  
\end{itemize}

\subsection{Writing \textit{Sütterlin} with \textit{Suetterlin-HJZ-Italic\_1911.ttf}}

{\suetterhjz Hier ein Text, geschrieben in Sütterlin via \texttt{xelatex} und Font \texttt{Suetterlin-HJZ-Italic\_1911.ttf}, entwickelt von Hans {Jürgen Zinken}}

\begin{itemize}
  \item \textit{Rundes-S}:
  \begin{itemize}
    \item[\faXmark] am Wortende: Haus = {\suetterhjz Hau\-s}, Maus = {\suetterhjz Maus}, Tennis = {\suetterhjz Tennis} 
    \item[\faXmark] in zusammengesetzten Wörtern: Marsmensch = {\suetterhjz Mars\-mensch}, Haustür = {\suetterhjz Haus\-tür}
    \item[\faXmark] bei Fugenmorphem: Königstochter = {\suetterhjz Königs\-tochter}, Valentins\-tag ={\suetterhjz Valentins\-tag}
    \item[\faXmark] vor konsonantischen Nachsilben: böse = {\suetterhjz böse} versus Bosheit = {\suetterhjz Bos\-heit}
    \item[\faXmark] am Ende der Vorsilben 'dis' und 'des': Desinformation = {\suetterhjz Des\-information}, Disharmonie = {\suetterhjz Dis\-harmonie}
    \item[\faXmark] am Wortende vor k: Kiosk = {\suetterhjz Kios\-k} grotesk = {\suetterhjz grotes\-k}
    \item[\faXmark] am Silbenende vor d, k, m, n, w: Potsdam = {\suetterhjz Pots\-dam}, Charisma= {\suetterhjz Charis\-ma}
  \end{itemize}
  \item[\faCheck] \textit{Langes s} im Wort und sonst: Mensch = {\suetterhjz Mensch}, hasten = {\suetterhjz hasten}, Tasse = {\suetterhjz Tasse}
  \item[\faCheck] ß wie früher üblich: Haß = {\suetterhjz Ha\ss} , daß = {\suetterhjz da\ss}
\end{itemize}

\section{Zur Verifikation von klassischen Zitierformen}


\begin{itemize}

  \item \enquote{Zitat mit \enquote{eingebettetem Zitat}}

  \item Deutscher Satz \foreignquote{german}{with embedded foreign phrase}
  als eingebettetes Zitat. Erwartetes Resultat: Deutsche Anführungszeichen,
  weil Teil in einem deutsche Satz.

  \item Autonomes englischsprachiges Zitat:
  \begin{quote}
    \foreignquote{english}{This shall be an English written paragraph containing
    a set of sentences which together build the quote.}
  \end{quote}
  Erwartetes Ergebnis: Englische Anführungszeichen, weil autonomer Satz.
\end{itemize}

\subsection{Verifikation: Literaturnachweis}
\begin{itemize}
  \item Buch1 (erstverwendungen)\footcite[vgl.][15]{KantKdrV1974}
  \item ders. ebda.\footcite[vgl.][15]{KantKdrV1974}
  \item ders. a.a.O.\footcite[vgl.][23]{KantKdrV1974}
  \item Buch2 (erstverwendung)\footcite[vgl.][15]{KantKdU1974}
  \item Buch1 (wiederverwendung)\footcite[vgl.][15]{KantKdrV1974}
\end{itemize}

% This file is part of proScientia.ltx
% (c) 2022 Karsten Reincke (https://github.com/kreincke/proScientia.ltx)
% It is distributed under the terms of the creative commons license
% CC-BY-4.0 (= https://creativecommons.org/licenses/by/4.0/)


% specific abbreviations
\abbr[utb]{UTB}{Uni-Taschenbuch}
\abbr[stw]{stw}{suhrkamp taschenbuch wissenschaft}

% general abbreviations
\abbr[vgl]{vgl.}{vergleiche}
\abbr[aaO]{a.a.O.}{am angegebenen Ort}
\abbr[ebda]{ebda.}{ebenda}
\abbr[ifross]{ifross}{Institut für Rechtsfragen der Freien und Open Source Software}
\abbr[wp]{wp.}{webpage = Webdokument ohne innere Seitennummerierung}
\abbr[RDL]{RDL}{Referenzdownload am}


% This file is part of proScientia.ltx
% (c) 2022 Karsten Reincke (https://github.com/kreincke/proScientia.ltx)
% It is distributed under the terms of the creative commons license
% CC-BY-4.0 (= https://creativecommons.org/licenses/by/4.0/)

\abbr[afda]{AfdA}{Anzeiger für deutsches Altertum}
\abbr[zfda]{ZfdA}{Zeitschrift für deutsches Altertum und deutsche Literatur [ISSN: 00442518]}
\abbr[zfaw]{}{Zeitschrift für Allgemeine Wissenschaftstheorie / Journal for General Philosophy of Science [ISSN: 0044-2216]}

\printnomenclature
\printbibliography

\end{document}
